\par\Needspace{12\baselineskip}\pdfbookmark[1]{Vega. Thesaurus Logarithmorum. 1851.}{vega1851}
\begin{center}
\rule{0.7\linewidth}{0.4pt}\par
Astronomische Nachrichten. Nr.\ 756. 1851. Mai 2.\par
\rule{0.7\linewidth}{0.4pt}
\end{center}
\begin{center}\emph{Einige Bemerkungen zu} \textsc{Vega}'s \emph{Thesaurus Logarithmorum, von Herrn}\par
\emph{Geheimen Hofrat} \textsc{Gauss}.\end{center}

\textsc{Vega}s \emph{Thesaurus Logarithmorum} ist bekanntlich größtenteils ein Neudruck der größeren \textsc{Vlacq}'schen Logarithmentafeln. In der Vorrede nennt \textsc{Vega} eine beträchtliche Anzahl von Fehlern im Original, die er verbessert hat. Er fügt hinzu, dass er außerdem eine sehr große Zahl von Unrichtigkeiten in der letzten Ziffer der Logarithmen berichtigt habe, die eine, zwei, drei oder vier Einheiten betrugen. Mit gleicher Sorgfalt seien auch die neu hinzugekommenen Tafeln berechnet, geprüft und berichtigt worden, insbesondere also die in \SourceLeaf{267} den beiden ersten Graden für jede einzelne Sekunde angegebenen Logarithmen der trigonometrischen Größen. \textsc{Vega} scheint sich nun mit der Hoffnung geschmeichelt zu haben, seine Tafeln seien auf diese Weise beinahe fehlerfrei geworden. Um \emph{vollkommen} fehlerfreie Tafeln zu erhalten, verspricht er für die erste Anzeige jedes etwa noch verbliebenen Fehlers, der eine falsche Rechnung verursachen könne (\emph{pro sphalmatibus calculum turbantibus}), eine Prämie von einem Dukaten zu bezahlen. Ob diese auf den $1^{\mathrm{sten}}$ Oktober $1794$ datierte Auslobung jemals Folgen hatte, ist mir nicht bekannt.

Ich bin unsicher, ob \textsc{Vega} sich ganz klargemacht hatte, welche Fehler als mögliche Ursache falscher Rechnungen zu betrachten seien. Für alle Tafeln, die theoretisch feststehende irrationale Größen darstellen sollen, gilt bekanntlich der Grundsatz, dass der Tabellenwert dem wahren Wert stets so nahe kommen soll, wie es bei der gewählten Zahl von Dezimalstellen möglich ist. Die Abweichung darf folglich niemals mehr als eine halbe Einheit der letzten Dezimalstelle betragen. Jeder Verstoß gegen diese strenge Norm ist ein Fehler, der einen Rechner, der sich auf die strenge Übereinstimmung verlässt, möglicherweise zu einem falschen Ergebnis führen kann. Gibt man diese strenge Auslegung auf und berücksichtigt bei seinem Urteil die \emph{Erheblichkeit} der Unrichtigkeit, so gerät die Entscheidung in den Bereich der Willkür. Der bereits erwähnte Umstand, dass \textsc{Vega} selbst von Korrekturen an den \textsc{Vlacq}'schen Tafeln spricht, die nur eine Einheit in der letzten Stelle betrugen, und dass er vollkommene Fehlerfreiheit als sein Ziel bezeichnet, scheint allerdings darauf hinzudeuten, dass er die strenge Auslegung im Sinn hatte. Auch habe ich den ersten Teil des Thesaurus, der die Logarithmen der fortlaufenden ganzen Zahlen enthält, bei sehr vielen gelegentlichen Vergleichen mit Bestimmungen auf mehr Stellen stets als sehr korrekt befunden.

Seit jener Zeit sind bei mehreren anderen Logarithmentafeln ähnliche Preise für die erste Anzeige von Fehlern in den Zahlen ausgesetzt worden, um ihre Richtigkeit zu vervollkommnen. Ich weiß jedoch nicht, ob diese Angebote Erfolg hatten, mit Ausnahme des $1847$ bei Tauchnitz in Leipzig von \textsc{Köhler} herausgegebenen logarithmisch-trigonometrischen Handbuchs. Der Verleger dieser Tafeln versprach beim ersten Erscheinen, für die erste \SourceLeaf{268} innerhalb einer festgesetzten Frist eingereichte Anzeige jedes Fehlers, der falsche Ergebnisse verursachen könne, einen Louisdor zu bezahlen. Nach einem gedruckten Bericht vom 1. Juli 1848 wurde diese Prämie für vier angezeigte Fehler tatsächlich ausgezahlt. Lobende Erwähnung verdient dabei, dass einer dieser Fehler jene Bedingung nur unter Anerkennung der obigen strengen Auslegung erfüllen konnte. Der Logarithmus von \(103000\), der auf \(12\) Stellen genau
\[
=5{,}012837224705
\]
ist, war nämlich in den ersten Drucken mit acht Ziffern als
\[
=5{,}01283723
\]
angegeben, während die nach dem Grundsatz gerundete Zahl \(5{,}01283722\) ist. Es ist zu wünschen, dass diese Entscheidung in ähnlichen künftigen Fällen als Präzedenzfall anerkannt wird.

Hätte \textsc{Vega} sich zu derselben strengen Auslegung bekannt, so hätte es ihm leicht wie König Shiram ergehen können, dessen Kornkammern nicht ausreichten, um dem Erfinder des Schachspiels die versprochene Belohnung zu gewähren.

\Needspace{42\baselineskip}
Dass die von \textsc{Vega} angebotene Belohnung ihn unter jener Voraussetzung teuer zu stehen kommen konnte, lässt sich schon ohne jedes Nachrechnen an einem Umstand erkennen, der leicht zu bestätigen ist, meines Wissens aber sonst noch nirgends zur Sprache gebracht wurde. In der Tafel der Logarithmen trigonometrischer Größen sind die Zahlen der Sinusspalte beinahe ausnahmslos\HistoricalNote{21} genau gleich der Summe der Zahlen der Kosinus- und der Tangentenspalte. Da alle diese Zahlen nur gerundete Werte der genauen irrationalen Größen sind, ist klar, dass jene Gleichheit bei streng richtiger Rundung in allen Fällen ausbleiben wird, in denen die Abweichungen von den genauen Werten in der zweiten und dritten Spalte dasselbe Vorzeichen haben und ihre Summe mehr als eine halbe Einheit der letzten Dezimalstelle beträgt. Man erkennt leicht, dass bei einer großen \SourceLeaf{269} Zahl von Fällen dieser Ausnahmefall im Durchschnitt einmal unter vieren vorkommen wird. Tatsächliche Zählungen in solchen siebenstelligen Tafeln, bei denen auf die Richtigkeit der letzten Ziffer sorgfältig geachtet wurde, bestätigen dies. (Eine solche Zählung in den erwähnten \textsc{Köhler}'schen Tafeln bei den \(900\) einzelnen Minuten von \(30^{0}0'\) bis \(44^{0}59'\) ergab beispielsweise genau \(225\) solche Fälle; diese vollkommene Übereinstimmung ist allerdings als zufällig anzusehen.) \textsc{Vega}s Thesaurus enthält die Logarithmen der Sinus, Kosinus und Tangenten von \(22680\) Bögen. Unter der Voraussetzung, dass sämtliche Zahlen in zwei dieser Spalten genau nach dem Grundsatz gerundet sind, darf man daher mit einer geringen Unsicherheit nach oben oder unten in der dritten Spalte \(5670\) Logarithmen erwarten, die um eine Einheit unrichtig angegeben sind. Für den Preis so vieler Dukaten wären wohl geeignete Personen zur vollständigen Neuberechnung bereit gewesen.

Welche beiden Spalten die ursprünglichen sind, lässt sich durch Vergleich mit anderen, auf mehr als zehn Stellen zuverlässigen Bestimmungen erkennen: schon an wenigen Bögen und mit Gewissheit, wenn die Voraussetzung der strengen Richtigkeit der ursprünglichen Spalten zutrifft; andernfalls aber wenigstens mit überwiegender Wahrscheinlichkeit an einer etwas größeren Zahl.

\Needspace{42\baselineskip}
Einiges dazu Brauchbare findet sich bereits am Schluss der Dezimaltafeln, die \textsc{Hobert} und \textsc{Ideler} 1799 veröffentlicht haben. Sie wurden zwar nur mit sieben Ziffern gedruckt, waren aber im Manuskript auf doppelt so viele ausgearbeitet. Dort werden \(138\) fehlerhafte Logarithmen aus \textsc{Vega}s Thesaurus angegeben. Um genau der Regel zu entsprechen, benötigen \(127\) davon eine Korrektur um eine Einheit in der zehnten Stelle, \(10\) eine Korrektur um zwei Einheiten und einer um drei Einheiten. Von den \(127\) entfallen \(49\) auf Sinus, \(35\) auf Kosinus und \(43\) auf Tangenten; die übrigen \(11\) Korrekturen um \(2\) oder \(3\) Einheiten betreffen ausschließlich Tangenten. Die betreffenden Bögen sind die durch \(27\) Minuten messbaren; die nicht mit aufgeführten sind diejenigen, deren Logarithmen keine Korrektur benötigten. Diese Vergleiche hätten übrigens auch ohne die handschriftlichen \textsc{Hobert-Ideler}'schen Tafeln vorgenommen werden können, da die allgemein verbreiteten \textsc{Callet}'schen Tafeln die Logarithmen der Sinus und Kosinus für alle Tausendstel des Quadranten auf \(14\) Ziffern enthalten. Entnimmt man ihnen außerdem das Benötigte für alle Bögen unter \SourceLeaf{270} zwei Grad, die durch \(5'24''\), aber nicht durch \(27'\) messbar sind, so erhält man einen Vergleich von \textsc{Vega}s Logarithmen für \(118\) Bögen mit genaueren Bestimmungen. Die Ergebnisse sind in der folgenden Übersicht (I) zusammengestellt:
\[
\begin{array}{rccc}
 & \mathrm{Sin.} & \mathrm{Cos.} & \mathrm{Tang.} \\
0 & 65 & 75 & 54 \\
1 & 53 & 43 & 53 \\
2 &    &    & 10 \\
3 &    &    & 1
\end{array}
\]
Die Zahlen etwa der letzten Spalte bedeuten, dass von \(118\) Tangentenlogarithmen \(54\) richtig angegeben sind, \(53\) eine Korrektur um eine Einheit in der zehnten Stelle benötigen, \(10\) eine Korrektur um zwei Einheiten und einer eine Korrektur um drei Einheiten. Daraus folgt bereits entschieden, dass \emph{keine} Spalte der \textsc{Vega}'schen Tafel durchgehend richtig angegeben ist, und mit überwiegender Wahrscheinlichkeit, dass die Tangentenlogarithmen die abgeleiteten sind: Sie entstehen durch einfache Subtraktion der Zahlen der Kosinusspalte von denen der Sinusspalte.

\Needspace{22\baselineskip}
Die Verfasser der erwähnten Dezimaltafeln hätten aus ihrem Manuskript übrigens doppelt so viele Vergleiche geben können. Eine noch viel größere Ausbeute erhält man jedoch aus der höchst wertvollen Tafel von \textsc{Brigg}'s in der nach dessen Tod von \textsc{Gellibrand} 1633 herausgegebenen Trigonometria Britannica. Sie liefert die Logarithmen der Sinus und Kosinus für alle Hundertstel der gewöhnlichen Grade auf \(14\) Ziffern, und \textsc{Callet} hat ihr die oben erwähnten Zahlen entnommen. Sie enthält das Material, um \textsc{Vega}s Tafeln bei \(1060\) Bögen, also insgesamt \(3180\) Logarithmen, prüfen zu können. Ich selbst habe diese Prüfung jedoch auf \(81\) Bögen oder \(243\) Logarithmen von \(14^{0}0'\) bis \(18^{0}0'\) beschränkt. Das Ergebnis folgt hier (II):
\[
\begin{array}{rccc}
 & \mathrm{Sin.} & \mathrm{Cos.} & \mathrm{Tang.} \\
0 & 29 & 56 & 36 \\
1 & 52 & 25 & 42 \\
2 &    &    & 3
\end{array}
\]
\par
\SourceLeaf{271}
Fasst man die Gruppen (I) und (II) zusammen und zählt dabei die beiden gemeinsamen Bögen nur einmal, so erhält man für 190 Bögen folgende Ausbeute (III):
\begin{center}
\begin{tabular}{rrrr}
 & Sin. & Cos. & Tang. \\
0 & 89 & 126 & 87 \\
1 & 101 & 64 & 90 \\
2 & & & 12 \\
3 & & & 1
\end{tabular}
\end{center}
Will man diese Zahlen zur Schätzung des Verhältnisses richtiger zu unrichtigen Logarithmen verwenden, muss man zunächst einen kleinen Abzug machen. In (I) ist der Logarithmus des Sinus von $45^{0}$ zweimal gezählt, nämlich zugleich auch als $\log\cos45^{0}$. Außerdem ist in der Tangentenspalte der Logarithmus der Tangente von $45^{0}$ mitgezählt, der jedoch rational ist. Dasselbe gilt für (III). Für den genannten Zweck muss man das Ergebnis daher so ausdrücken: Von 568 geprüften irrationalen Logarithmen erwiesen sich 301 als richtig und 267 als unrichtig. Dürfte man dieses Verhältnis als im Durchschnitt zutreffend ansehen, so wären unter den 68038 irrationalen Logarithmen in \textsc{Vega}s Thesaurus nach der Wahrscheinlichkeit etwa 31983 fehlerhafte anzunehmen; die Kotangenten sind dabei, wie angemessen, nicht mitgezählt.

\Needspace{36\baselineskip}
Wahrscheinlich ist diese Zahl aber noch erheblich zu klein. In meinen Papieren finde ich für einige Gruppen von Bögen eine vor längerer Zeit zu anderen Zwecken auf 14 Ziffern ausgeführte Berechnung der trigonometrischen Logarithmen. Diese Bögen schreiten nicht sprunghaft, sondern in denselben Abständen wie \textsc{Vega}s Tafeln fort. Daraus geht wenigstens hervor, dass das obige Ergebnis (III) noch keinen richtigen Maßstab für die Ungenauigkeit dieser Tafeln liefert. Während unter jenen 190 Bögen kein einziger Fall vorkommt, in dem der Logarithmus eines Sinus oder Kosinus um mehr als \textit{eine} Einheit verbessert werden müsste, sind solche Fälle bei den Bögen keineswegs selten, die \textit{nicht} in der Trigonometria Britannica vorkommen und für die das Vergleichsmaterial daher erst durch eine besondere Rechnung beschafft werden muss. Ich füge deshalb die Ergebnisse jener Rechnungen bei, da sie unserer Vorstellung vom Grad der Ungenauigkeit der Zahlen in \textsc{Vega}s Thesaurus eine festere Grundlage geben können.
\par
\SourceLeaf{272}
Der Vergleich der Zahlen im Thesaurus für die 21 Bögen von $15^{0}38'20''$ bis $15^{0}41'40''$ mit genauer berechneten Werten hat ergeben (IV):
\begin{center}\begin{tabular}{r r r r}
& Sin. & Cos. & Tang.\\
0&4&12&1\\
1&9&8&8\\
2&6&1&6\\
3&2&&4\\
4&&&2
\end{tabular}\end{center}
Die größten Abweichungen treten bei folgenden Bögen auf:
\[
\begin{array}{l l r r r}
15^{0}40'20'',&\text{wo die Korrekturen}&+3,&-1,&+4\\
15\quad41\quad30,&\text{wo folgende}&+3,&0,&+4
\end{array}
\]
an den Logarithmen des Sinus, des Kosinus und der Tangente anzubringen sind.

\Needspace{11\baselineskip}Ebenso ergab der Vergleich von \textsc{Vega}s Zahlen für die 93 Bögen von $1^{0}19'52''$ bis $1^{0}21'24''$ mit genauerer Rechnung Folgendes (V):
\begin{center}\begin{tabular}{r r r r}
& Sin. & Cos. & Tang.\\
0&38&30&17\\
1&39&56&41\\
2&15&7&22\\
3&1&&11\\
4&&&2
\end{tabular}\end{center}
Die größten Abweichungen treten bei den Bögen $1^{0}20'10''$ und $1^{0}20'15''$ auf; die nötigen Korrekturen betragen beim ersten $-2,+1,-4$ und beim anderen $-2,+2,-4$.

\Needspace{28\baselineskip}
Die Gruppen IV und V haben nicht denselben Ursprung. Die Zahlen des Thesaurus für Gruppe IV stammen aus \textsc{Vlacq}'s Trigonometria Artificialis; die anderen gehören zu den neu hinzugekommenen, die unter \textsc{Vega}s Leitung und Aufsicht von Leutnant \textsc{Dorfmund} berechnet wurden. Nach den obigen Übersichten erscheinen die letzteren etwas weniger ungenau als die ersteren, obwohl Gruppe IV zu wenige Fälle für ein sicheres Urteil umfasst. Jedenfalls liefert die Zusammenfassung beider Gruppen eine Schätzung \SourceLeaf{273} der gesamten Ungenauigkeit der Tafeln, die eher etwas zu günstig als umgekehrt ausfallen wird. Die Zusammenfassung ergibt für 114 Bögen (VI):
\begin{center}\begin{tabular}{r r r r}
& Sin. & Cos. & Tang.\\
0&42&42&18\\
1&48&64&49\\
2&21&8&28\\
3&3&&15\\
4&&&4
\end{tabular}\end{center}
Unter 342 Logarithmen sind hier also nur 102. die keiner Verbesserung bedürfen, gegenüber 240 ungenauen. Nach diesem Verhältnis wären unter den 68038 irrationalen Logarithmen 47746 ungenaue zu erwarten.

Die Summe der Quadrate der Abweichungen ergibt sich für die Sinus zu 159, für die Kosinus zu 96 und für die Tangenten zu 360. Als mittleren Fehler kann man daher für die Sinus $1{,}18$, für die Kosinus $0{,}92$ und für die Tangenten $1{,}78$ annehmen. Danach kann man nicht daran zweifeln, dass die Tangentenlogarithmen die abgeleiteten sind.

\begin{samepage}
Dass die Zahlen der Kosinusspalte weniger ungenau sind als die der Sinusspalte, liegt wohl ohne Zweifel wenigstens teilweise daran, dass bei den ersteren die zum Auffüllen erforderlichen Interpolationsmethoden einfacher sind. Möglicherweise haben jedoch noch andere Ursachen mitgewirkt, über die sich nur unsichere Vermutungen aufstellen ließen.
\begin{center}\bigskip\rule{.3\linewidth}{.4pt}\end{center}
\end{samepage}
